\begin{table}[htbp]
\centering
\caption{Sample composition, by estimation sample}
\label{tab:sample_composition}
\begin{tabular}{lcc}
\toprule
 & Revenue per capita & Revenue per store / density \\
\midrule
County-year observations & 1,283 & 964 \\
Unique counties & 178 & 165 \\
\addlinespace
Border counties (ever, in sample) & 54 & 48 \\
Interior counties (ever, in sample) & 138 & 128 \\
Counties that switch classification & 14 & 11 \\
\addlinespace
Border county-years & 280 & 219 \\
Interior county-years & 1,003 & 745 \\
\addlinespace
Capped-entry county-years & 1,024 & 840 \\
Free-entry county-years & 259 & 124 \\
\bottomrule
\end{tabular}
\begin{minipage}{0.92\textwidth}
\vspace{0.3em}
\footnotesize
\textit{Notes:} Both columns exclude New Mexico's incomplete 2026 partial year and match
Table~\ref{tab:main_pooled} exactly. ``Border (interior) counties, ever'' counts a county once if it is
ever classified as border (interior) in at least one sample year; because the border classification is
time-varying (Section~\ref{sec:timevarying}), a county can appear in both categories across different
years -- ``counties that switch classification'' reports how many do. Capped/free-entry county-years use
the state-year \texttt{Capped} indicator (Section~\ref{sec:institutional}).
\end{minipage}
\end{table}
